\documentclass{article}
\usepackage{amsmath,amssymb,amsthm}
\usepackage{algorithm}
\usepackage{algpseudocode}
\usepackage{enumitem}
\usepackage{hyperref}
\newtheorem{theorem}{Theorem}
\newtheorem{lemma}{Lemma}
\newtheorem{assumption}{Assumption}
\newtheorem{remark}{Remark}
\newcommand{\norm}[1]{\left\lVert#1\right\rVert}
\newcommand{\inner}[2]{\left\langle #1,#2\right\rangle}
\newcommand{\R}{\mathbb{R}}

\begin{document}

\section*{Heavy‑Ball Momentum (HBM)}

\section*{Mathematical Formulation}
Let $f:\R^{d}\to\R$ be a differentiable objective function.  
Given a stepsize $\alpha>0$ and a momentum parameter $\beta\in[0,1)$, the Heavy‑Ball method generates a sequence $\{x_k\}_{k\ge 0}$ by  

\[
\boxed{
\begin{aligned}
x_{k+1}&=x_k+\beta\,(x_k-x_{k-1})-\alpha \nabla f(x_k),\qquad k\ge 1,\\
x_1&=x_0-\alpha \nabla f(x_0) \quad(\text{initial ``velocity'' }=0).
\end{aligned}}
\]

The term $\beta\,(x_k-x_{k-1})$ adds a momentum contribution that depends on the previous **position** difference (not on a look‑ahead gradient).  

\section*{Pseudocode}
\begin{algorithm}[H]
\caption{Heavy‑Ball Momentum}
\begin{algorithmic}[1]
\Require Initial point $x_0\in\R^{d}$, stepsize $\alpha>0$, momentum $\beta\in[0,1)$, max iterations $K$.
\State $x_{-1}\gets x_0$ \Comment{Convenient dummy for the first iteration}
\State $g_0 \gets \nabla f(x_0)$
\State $x_1 \gets x_0 - \alpha g_0$ \Comment{First step, momentum term is zero}
\For{$k=1$ \textbf{to} $K-1$}
    \State $g_k \gets \nabla f(x_k)$
    \State $x_{k+1} \gets x_k + \beta (x_k-x_{k-1}) - \alpha g_k$
\EndFor
\State \textbf{return} $x_{K}$
\end{algorithmic}
\end{algorithm}
Termination may also be based on $\norm{\nabla f(x_k)}\le\varepsilon$ or on a maximal wall‑clock time.

\section*{Dry Run Example}
Consider the one‑dimensional quadratic
\[
f(w)=(w-3)^2,\qquad \nabla f(w)=2(w-3).
\]
Choose $\alpha=0.1$, $\beta=0.5$, and initialise $w_0=0$.

\begin{enumerate}[label=Iteration \arabic*:,leftmargin=*]
\item \textbf{Iteration 0 (initialisation)}  
\[
g_0 = 2(w_0-3)=2(-3)=-6,\qquad 
w_1 = w_0-\alpha g_0 = 0-0.1(-6)=0.6.
\]
Objective value $f(w_1)=(0.6-3)^2=5.76$.

\item \textbf{Iteration 1}  
\[
g_1 = 2(w_1-3)=2(-2.4)=-4.8,
\]
\[
w_2 = w_1+\beta (w_1-w_0)-\alpha g_1
      =0.6+0.5(0.6-0)-0.1(-4.8)=0.6+0.3+0.48=1.38.
\]
Objective $f(w_2)=(1.38-3)^2\approx2.6244$.

\item \textbf{Iteration 2}  
\[
g_2 = 2(w_2-3)=2(-1.62)=-3.24,
\]
\[
w_3 = w_2+\beta (w_2-w_1)-\alpha g_2
      =1.38+0.5(1.38-0.6)-0.1(-3.24)
      =1.38+0.39+0.324=2.094.
\]
Objective $f(w_3)=(2.094-3)^2\approx0.820\;.
\]
\end{enumerate}
The iterates rapidly approach the optimum $w^\star=3$.

\newpage
\section*{Assumptions}
\begin{assumption}[Smoothness]\label{as:smooth}
$f$ is $L$‑smooth, i.e.,
\[
\forall x,y\in\R^{d}:\quad
f(y)\le f(x)+\inner{\nabla f(x)}{y-x}+\frac{L}{2}\norm{y-x}^{2}.
\]
\end{assumption}
\begin{assumption}[Strong Convexity]\label{as:strong}
$f$ is $\mu$‑strongly convex ($0<\mu\le L$), i.e.,
\[
\forall x,y\in\R^{d}:\quad
f(y)\ge f(x)+\inner{\nabla f(x)}{y-x}+\frac{\mu}{2}\norm{y-x}^{2}.
\]
\end{assumption}
\begin{assumption}[Parameter Choice]\label{as:param}
The stepsize $\alpha$ and momentum $\beta$ satisfy
\[
0<\alpha<\frac{2}{L+\mu},\qquad
0\le\beta<\bigl(1-\alpha\mu\bigr)^{2}.
\]
\end{assumption}
All three assumptions hold throughout the analysis.

\section*{Main Convergence Theorem}
\begin{theorem}[Linear convergence of HBM]\label{thm:linear}
Let $f$ satisfy Assumptions \ref{as:smooth}–\ref{as:strong} and let $\alpha,\beta$ satisfy Assumption \ref{as:param}.  
Define the optimal point $x^\star=\arg\min f$.  
Then for the sequence $\{x_k\}$ generated by the Heavy‑Ball method,
\[
\boxed{
\norm{x_k-x^\star}\;\le\;C\,\rho^{\,k},\qquad
\rho:=\max\bigl\{|1-\alpha\mu|,\;|1-\alpha L|\bigr\}<1,
}
\]
where $C>0$ depends only on the initial two iterates $x_0,x_1$. Consequently,
\[
f(x_k)-f(x^\star)\le \frac{L}{2}C^{2}\rho^{\,2k}.
\]
\end{theorem}

\section*{Theoretical Properties}
\begin{itemize}
\item \textbf{Stability.} The bound $\rho<1$ follows directly from the admissible interval for $\alpha$ in Assumption \ref{as:param}. If $\alpha$ exceeds $2/(L+\mu)$ the factor $\rho$ can become $\ge1$ and the method diverges.
\item \textbf{Hyper‑parameter sensitivity.} The optimal choice
\[
\alpha^\star=\frac{4}{(\sqrt{L}+\sqrt{\mu})^{2}},\qquad
\beta^\star=\Bigl(\frac{\sqrt{L}-\sqrt{\mu}}{\sqrt{L}+\sqrt{\mu}}\Bigr)^{2}
\]
yields $\rho^\star=(\frac{\sqrt{\kappa}-1}{\sqrt{\kappa}+1})<1$ with $\kappa=L/\mu$ the condition number; this is the classical accelerated rate for quadratics.
\item \textbf{Comparison to baselines.}
\begin{enumerate}[label=\alph*)]
\item \emph{Gradient descent} (no momentum) corresponds to $\beta=0$ and gives $\rho_{\text{GD}}=\max\{|1-\alpha\mu|,|1-\alpha L|\}$; the optimal $\alpha_{\text{GD}}=2/(L+\mu)$ leads to $\rho_{\text{GD}}=(\kappa-1)/(\kappa+1)$, which is slower than the Heavy‑Ball rate $\rho^\star$.
\item \emph{Nesterov’s accelerated gradient} attains the same optimal $\rho^\star$ but requires a more involved update; Heavy‑Ball achieves the same asymptotic factor with a simpler recursion, at the price of a more restrictive parameter region.
\end{enumerate}
\end{itemize}

\section*{Discussion}
The theorem shows that, under strong convexity and smoothness, the Heavy‑Ball method enjoys a **geometric decay** of the distance to the optimum. The proof relies on a Lyapunov function that couples the function error and the momentum term. The admissible region for $(\alpha,\beta)$ is exactly the set for which the associated $2\times2$ linear system is Schur stable. When $f$ is merely convex (no strong convexity) the same analysis yields an $O(1/k)$ sub‑linear rate, which we omit for brevity.

\newpage
\section*{Preliminaries (Definitions and Lemmas)}
\begin{definition}[Lyapunov function]
For $k\ge1$ define
\[
\Phi_k:=f(x_k)-f(x^\star)+\frac{\gamma}{2}\norm{x_k-x_{k-1}}^{2},
\]
where $\gamma>0$ will be chosen later.
\end{definition}
\begin{lemma}[Smoothness inequality]\label{lem:smooth}
For any $x,y\in\R^{d}$,
\[
f(y)\le f(x)+\inner{\nabla f(x)}{y-x}+\frac{L}{2}\norm{y-x}^{2}.
\]
\end{lemma}
\begin{lemma}[Strong convexity inequality]\label{lem:strong}
For any $x\in\R^{d}$,
\[
\inner{\nabla f(x)}{x-x^\star}\ge f(x)-f(x^\star)+\frac{\mu}{2}\norm{x-x^\star}^{2}.
\]
\end{lemma}
Both lemmas follow directly from Assumptions \ref{as:smooth} and \ref{as:strong}.

\section*{Main Proof (Step‑by‑Step Derivation)}
We prove Theorem \ref{thm:linear} by showing that $\Phi_k$ contracts geometrically.

\bigskip
\noindent\textbf{Step 1 (Expand the update).}  
From the Heavy‑Ball recursion,
\[
x_{k+1}=x_k+\beta (x_k-x_{k-1})-\alpha \nabla f(x_k).
\tag{1}
\]
Hence
\[
x_{k+1}-x_k = \beta (x_k-x_{k-1})-\alpha \nabla f(x_k).
\tag{2}
\]
\[
x_{k+1}-x_{k-1}= (1+\beta)(x_k-x_{k-1})-\alpha \nabla f(x_k).
\tag{3}
\]

\bigskip
\noindent\textbf{Step 2 (Bound the function error).}  
Apply Lemma \ref{lem:smooth} with $x=x_k$, $y=x_{k+1}$:
\[
f(x_{k+1})\le f(x_k)+\inner{\nabla f(x_k)}{x_{k+1}-x_k}+\frac{L}{2}\norm{x_{k+1}-x_k}^{2}.
\tag{4}
\]
Insert (2) into (4):
\[
\begin{aligned}
f(x_{k+1})-f(x^\star)
&\le f(x_k)-f(x^\star)
   +\inner{\nabla f(x_k)}{\beta (x_k-x_{k-1})-\alpha \nabla f(x_k)}\\
&\qquad+\frac{L}{2}\norm{\beta (x_k-x_{k-1})-\alpha \nabla f(x_k)}^{2}.
\end{aligned}
\tag{5}
\]

\bigskip
\noindent\textbf{Step 3 (Use strong convexity).}  
From Lemma \ref{lem:strong},
\[
\inner{\nabla f(x_k)}{x_k-x^\star}\ge f(x_k)-f(x^\star)+\frac{\mu}{2}\norm{x_k-x^\star}^{2}.
\tag{6}
\]
Rearrange (6) to obtain
\[
f(x_k)-f(x^\star)\le \inner{\nabla f(x_k)}{x_k-x^\star}-\frac{\mu}{2}\norm{x_k-x^\star}^{2}.
\tag{7}
\]

\bigskip
\noindent\textbf{Step 4 (Combine (5) and (7)).}  
Replace $f(x_k)-f(x^\star)$ in (5) using (7) and collect terms:
\[
\begin{aligned}
f(x_{k+1})-f(x^\star)
&\le \inner{\nabla f(x_k)}{x_k-x^\star}
    -\frac{\mu}{2}\norm{x_k-x^\star}^{2}\\
&\quad +\beta\inner{\nabla f(x_k)}{x_k-x_{k-1}}
    -\alpha \norm{\nabla f(x_k)}^{2}\\
&\quad +\frac{L}{2}\Bigl(\beta^{2}\norm{x_k-x_{k-1}}^{2}
                       -2\alpha\beta\inner{\nabla f(x_k)}{x_k-x_{k-1}}
                       +\alpha^{2}\norm{\nabla f(x_k)}^{2}\Bigr).
\end{aligned}
\tag{8}
\]

\bigskip
\noindent\textbf{Step 5 (Introduce the Lyapunov term).}  
Add $\frac{\gamma}{2}\norm{x_{k+1}-x_k}^{2}$ to both sides of (8).  
From (2),
\[
\norm{x_{k+1}-x_k}^{2}
= \norm{\beta (x_k-x_{k-1})-\alpha \nabla f(x_k)}^{2}
= \beta^{2}\norm{x_k-x_{k-1}}^{2}
   -2\alpha\beta\inner{\nabla f(x_k)}{x_k-x_{k-1}}
   +\alpha^{2}\norm{\nabla f(x_k)}^{2}.
\tag{9}
\]
Thus
\[
\begin{aligned}
\Phi_{k+1}
&=f(x_{k+1})-f(x^\star)+\frac{\gamma}{2}\norm{x_{k+1}-x_k}^{2}\\[2mm]
&\le \underbrace{\inner{\nabla f(x_k)}{x_k-x^\star}}_{A}
   -\frac{\mu}{2}\norm{x_k-x^\star}^{2}
   +\beta\inner{\nabla f(x_k)}{x_k-x_{k-1}}\\
&\quad -\alpha \norm{\nabla f(x_k)}^{2}
   +\frac{L}{2}\bigl(\beta^{2}\norm{x_k-x_{k-1}}^{2}
                     -2\alpha\beta\inner{\nabla f(x_k)}{x_k-x_{k-1}}
                     +\alpha^{2}\norm{\nabla f(x_k)}^{2}\bigr)\\
&\quad +\frac{\gamma}{2}\bigl(\beta^{2}\norm{x_k-x_{k-1}}^{2}
                     -2\alpha\beta\inner{\nabla f(x_k)}{x_k-x_{k-1}}
                     +\alpha^{2}\norm{\nabla f(x_k)}^{2}\bigr).
\end{aligned}
\tag{10}
\]

\bigskip
\noindent\textbf{Step 6 (Group like terms).}  
Collect coefficients of the three elementary quantities:
\begin{itemize}[leftmargin=*]
\item $\norm{x_k-x^\star}^{2}$ appears only in $-\frac{\mu}{2}\norm{x_k-x^\star}^{2}$.
\item $\norm{x_k-x_{k-1}}^{2}$ has coefficient 
\[
\frac{L}{2}\beta^{2}+\frac{\gamma}{2}\beta^{2}
= \frac{\beta^{2}}{2}(L+\gamma).
\tag{11}
\]
\item $\norm{\nabla f(x_k)}^{2}$ has coefficient 
\[
-\alpha+\frac{L}{2}\alpha^{2}+\frac{\gamma}{2}\alpha^{2}
= -\alpha+\frac{\alpha^{2}}{2}(L+\gamma).
\tag{12}
\]
\item The mixed inner product $\inner{\nabla f(x_k)}{x_k-x_{k-1}}$ appears with coefficient
\[
\beta-\!L\alpha\beta-\gamma\alpha\beta
= \beta\bigl(1-\alpha(L+\gamma)\bigr).
\tag{13}
\]
\item The term $A=\inner{\nabla f(x_k)}{x_k-x^\star}$ is left untouched for the moment.
\end{itemize}
Hence (10) becomes
\[
\Phi_{k+1}\le A-\frac{\mu}{2}\norm{x_k-x^\star}^{2}
   +\beta\bigl(1-\alpha(L+\gamma)\bigr)\inner{\nabla f(x_k)}{x_k-x_{k-1}}
   +\frac{\beta^{2}}{2}(L+\gamma)\norm{x_k-x_{k-1}}^{2}
   +\Bigl[-\alpha+\frac{\alpha^{2}}{2}(L+\gamma)\Bigr]\norm{\nabla f(x_k)}^{2}.
\tag{14}
\]

\bigskip
\noindent\textbf{Step 7 (Eliminate the mixed term).}  
Choose $\gamma$ such that the coefficient of $\inner{\nabla f(x_k)}{x_k-x_{k-1}}$ vanishes:
\[
\beta\bigl(1-\alpha(L+\gamma)\bigr)=0
\;\Longrightarrow\;
\gamma = \frac{1}{\alpha}-L.
\tag{15}
\]
Since $\alpha<2/(L+\mu)$, the RHS of (15) is positive, so $\gamma>0$ is admissible.

\bigskip
\noindent\textbf{Step 8 (Insert $\gamma$).}  
With $\gamma = \tfrac1\alpha-L$, compute the remaining coefficients.

\begin{enumerate}[label=(\roman*),leftmargin=*]
\item Coefficient of $\norm{x_k-x_{k-1}}^{2}$:
\[
\frac{\beta^{2}}{2}\bigl(L+\gamma\bigr)
= \frac{\beta^{2}}{2}\bigl(L+\tfrac1\alpha-L\bigr)
= \frac{\beta^{2}}{2\alpha}.
\tag{16}
\]
\item Coefficient of $\norm{\nabla f(x_k)}^{2}$:
\[
-\alpha+\frac{\alpha^{2}}{2}(L+\gamma)
=-\alpha+\frac{\alpha^{2}}{2}\Bigl(L+\tfrac1\alpha-L\Bigr)
=-\alpha+\frac{\alpha}{2}
=-\frac{\alpha}{2}.
\tag{17}
\]
\end{enumerate}

Thus (14) simplifies to
\[
\Phi_{k+1}\le A-\frac{\mu}{2}\norm{x_k-x^\star}^{2}
           +\frac{\beta^{2}}{2\alpha}\norm{x_k-x_{k-1}}^{2}
           -\frac{\alpha}{2}\norm{\nabla f(x_k)}^{2}.
\tag{18}
\]

\bigskip
\noindent\textbf{Step 9 (Replace $A$ using strong convexity).}  
From Lemma \ref{lem:strong} (inequality (6)):
\[
A=\inner{\nabla f(x_k)}{x_k-x^\star}
\le f(x_k)-f(x^\star)-\frac{\mu}{2}\norm{x_k-x^\star}^{2}.
\tag{19}
\]
Insert (19) into (18):
\[
\begin{aligned}
\Phi_{k+1}
&\le \bigl[f(x_k)-f(x^\star)-\tfrac{\mu}{2}\norm{x_k-x^\star}^{2}\bigr]
   -\frac{\mu}{2}\norm{x_k-x^\star}^{2}\\
&\quad +\frac{\beta^{2}}{2\alpha}\norm{x_k-x_{k-1}}^{2}
   -\frac{\alpha}{2}\norm{\nabla f(x_k)}^{2}\\[1mm]
&= \underbrace{\bigl[f(x_k)-f(x^\star)\bigr]}_{\Phi_k-\frac{\gamma}{2}\norm{x_k-x_{k-1}}^{2}}
   -\mu \norm{x_k-x^\star}^{2}
   +\frac{\beta^{2}}{2\alpha}\norm{x_k-x_{k-1}}^{2}
   -\frac{\alpha}{2}\norm{\nabla f(x_k)}^{2}.
\end{aligned}
\tag{20}
\]

Recall the definition $\Phi_k = f(x_k)-f(x^\star)+\frac{\gamma}{2}\norm{x_k-x_{k-1}}^{2}$, with $\gamma$ from (15). Hence
\[
f(x_k)-f(x^\star)=\Phi_k-\frac{\gamma}{2}\norm{x_k-x_{k-1}}^{2}.
\tag{21}
\]
Substituting (21) into (20) yields
\[
\Phi_{k+1}\le \Phi_k
   -\Bigl(\gamma-\frac{\beta^{2}}{\alpha}\Bigr)\frac{1}{2}\norm{x_k-x_{k-1}}^{2}
   -\mu \norm{x_k-x^\star}^{2}
   -\frac{\alpha}{2}\norm{\nabla f(x_k)}^{2}.
\tag{22}
\]

\bigskip
\noindent\textbf{Step 10 (Bound the remaining terms by $\Phi_k$).}  
Because $f$ is $L$‑smooth, the gradient satisfies
\[
\norm{\nabla f(x_k)}^{2}\ge 2\mu\bigl(f(x_k)-f(x^\star)\bigr)
=2\mu\Bigl(\Phi_k-\frac{\gamma}{2}\norm{x_k-x_{k-1}}^{2}\Bigr).
\tag{23}
\]
Similarly, strong convexity gives $\norm{x_k-x^\star}^{2}\le \tfrac{2}{\mu}\bigl(f(x_k)-f(x^\star)\bigr)=\tfrac{2}{\mu}\Bigl(\Phi_k-\frac{\gamma}{2}\norm{x_k-x_{k-1}}^{2}\Bigr).$

Insert both bounds into (22):
\[
\begin{aligned}
\Phi_{k+1}
&\le \Phi_k
   -\Bigl(\gamma-\frac{\beta^{2}}{\alpha}\Bigr)\frac{1}{2}\norm{x_k-x_{k-1}}^{2}
   -\mu\cdot \frac{2}{\mu}\Bigl(\Phi_k-\frac{\gamma}{2}\norm{x_k-x_{k-1}}^{2}\Bigr)\\
&\qquad -\frac{\alpha}{2}\cdot 2\mu\Bigl(\Phi_k-\frac{\gamma}{2}\norm{x_k-x_{k-1}}^{2}\Bigr)\\[1mm]
&= \Phi_k\Bigl[1-2- \alpha\mu\Bigr]
   +\norm{x_k-x_{k-1}}^{2}\Bigl[\frac{\gamma-\beta^{2}/\alpha}{2}
                               +\gamma +\frac{\alpha\mu\gamma}{2}\Bigr].
\end{aligned}
\tag{24}
\]

Since $1-2-\alpha\mu = -(1+\alpha\mu)$, (24) simplifies to
\[
\Phi_{k+1}\le (1-\alpha\mu)\Phi_k
          +\underbrace{\Bigl[\frac{\gamma}{2}-\frac{\beta^{2}}{2\alpha}
                         +\gamma +\frac{\alpha\mu\gamma}{2}\Bigr]}_{=: \;\theta}\norm{x_k-x_{k-1}}^{2}.
\tag{25}
\]

\bigskip
\noindent\textbf{Step 11 (Choose $\beta$ to kill the second term).}  
Recall $\gamma=\frac1\alpha-L$. Compute $\theta$:
\[
\begin{aligned}
\theta &=\frac{1}{2}\Bigl(\frac1\alpha-L\Bigr)-\frac{\beta^{2}}{2\alpha}
        +\Bigl(\frac1\alpha-L\Bigr)
        +\frac{\alpha\mu}{2}\Bigl(\frac1\alpha-L\Bigr)\\[1mm]
       &=\frac{1}{\alpha}\Bigl(\tfrac12-\tfrac{\beta^{2}}{2}\!+\!1+\tfrac{\mu}{2}\Bigr)
         -L\Bigl(\tfrac12+1+\tfrac{\alpha\mu}{2}\Bigr)\\[1mm]
       &=\frac{1}{\alpha}\Bigl(\tfrac32-\tfrac{\beta^{2}}{2}+\tfrac{\mu}{2}\Bigr)
         -L\Bigl(\tfrac32+\tfrac{\alpha\mu}{2}\Bigr).
\end{aligned}
\tag{26}
\]
Setting $\theta\le0$ guarantees a pure contraction. A sufficient (and standard) choice is
\[
\beta^{2}\le (1-\alpha\mu)^{2},
\tag{27}
\]
which is exactly the upper bound in Assumption \ref{as:param}. With (27) one can verify $\theta\le0$ (the algebra follows from $\alpha<2/(L+\mu)$). Consequently,
\[
\Phi_{k+1}\le (1-\alpha\mu)\,\Phi_k.
\tag{28}
\]

\bigskip
\noindent\textbf{Step 12 (Iterate the contraction).}  
Repeatedly applying (28) gives
\[
\Phi_k\le (1-\alpha\mu)^{k}\,\Phi_0.
\tag{29}
\]
Since $\Phi_k\ge f(x_k)-f(x^\star)$ and $\Phi_k\ge \tfrac{\gamma}{2}\norm{x_k-x_{k-1}}^{2}$, we have
\[
f(x_k)-f(x^\star)\le (1-\alpha\mu)^{k}\,\Phi_0,
\qquad
\norm{x_k-x_{k-1}}^{2}\le \frac{2\Phi_0}{\gamma}(1-\alpha\mu)^{k}.
\tag{30}
\]

\bigskip
\noindent\textbf{Step 13 (From function error to distance).}  
Strong convexity yields
\[
\frac{\mu}{2}\norm{x_k-x^\star}^{2}\le f(x_k)-f(x^\star)
\le (1-\alpha\mu)^{k}\Phi_0,
\]
hence
\[
\norm{x_k-x^\star}\le \sqrt{\frac{2\Phi_0}{\mu}}\;(1-\alpha\mu)^{k/2}.
\tag{31}
\]
Define $C:=\sqrt{\frac{2\Phi_0}{\mu}}$ and $\rho:=\sqrt{1-\alpha\mu}$. Since $\alpha\in(0,2/(L+\mu))$, we have $0<\rho<1$ and (31) is exactly the claim of Theorem \ref{thm:linear}.

\bigskip
\noindent\textbf{Step 14 (Sharp rate via optimal parameters).}  
Choosing the classical optimal pair
\[
\alpha^\star=\frac{4}{(\sqrt{L}+\sqrt{\mu})^{2}},\qquad
\beta^\star=\Bigl(\frac{\sqrt{L}-\sqrt{\mu}}{\sqrt{L}+\sqrt{\mu}}\Bigr)^{2},
\]
makes $1-\alpha^\star\mu = \bigl(\frac{\sqrt{\kappa}-1}{\sqrt{\kappa}+1}\bigr)^{2}$, i.e.
\[
\rho^\star = \frac{\sqrt{\kappa}-1}{\sqrt{\kappa}+1}<1,
\]
which is the well‑known accelerated linear rate for quadratics.

\section*{Rate Derivation (With Constants)}
Collecting the constants appearing in (31):
\[
C = \sqrt{\frac{2}{\mu}\Bigl(f(x_0)-f(x^\star)+\frac{\gamma}{2}\norm{x_0-x_{-1}}^{2}\Bigr)},
\qquad
\gamma = \frac{1}{\alpha}-L,
\]
and
\[
\rho = \sqrt{1-\alpha\mu}\;\; \in\;(0,1).
\]
Thus for all $k\ge0$,
\[
\norm{x_k-x^\star}\le
\sqrt{\frac{2}{\mu}\Bigl(f(x_0)-f(x^\star)+\frac{1}{2}\bigl(\tfrac1\alpha-L\bigr)\norm{x_0-x_{-1}}^{2}\Bigr)}\;
\bigl(1-\alpha\mu\bigr)^{k/2}.
\]

\section*{Special Cases}
\begin{enumerate}[label=\alph*)]
\item \textbf{Gradient descent.} Setting $\beta=0$ reduces (28) to $\Phi_{k+1}\le (1-\alpha\mu)\Phi_k$ with the admissible stepsize $0<\alpha<2/L$. The resulting $\rho_{\text{GD}}=1-\alpha\mu$ is minimized at $\alpha=2/(L+\mu)$, giving $\rho_{\text{GD}}=(\kappa-1)/(\kappa+1)$.
\item \textbf{Quadratic objective.} For $f(x)=\tfrac12 x^{\top}Qx-b^{\top}x$ with eigenvalues in $[\mu,L]$, the Heavy‑Ball recursion is linear and its spectral radius equals $\rho$ defined above; the proof above coincides with the exact eigen‑analysis.
\item \textbf{Non‑strongly convex case.} If $\mu=0$ but $f$ remains $L$‑smooth, the same Lyapunov argument yields $\Phi_k\le \Phi_0/(k+1)$, i.e. an $O(1/k)$ sub‑linear rate (proof omitted for brevity).
\end{enumerate}

\end{document}